% TREŚĆ ROZDZIAŁU: zastąp wskazówki opisem testów i uzyskanych wyników.
% Każdą tabelę i każdy rysunek omów w tekście. Podawaj jednostki, warunki
% pomiaru oraz odwołanie do wymagania, które dany test weryfikuje.
\chapter{Weryfikacja i wyniki}
\label{chap:wyniki}

\section{Metodyka testów}
Opisz stanowisko, dane, warunki, procedurę i miary. Informacje powinny wystarczyć do powtórzenia najważniejszych testów.

\section{Wyniki}
Dane porównawcze obu algorytmów zawiera \cref{tab:porownanie-algorytmow}.
Zależność czasu sortowania bąbelkowego od liczby elementów pokazuje \cref{fig:wykres-sortowania}.
Kod funkcji wyszukującej maksimum zamieszczono w \cref{lst:max-element}.
\subsection{Porównanie algorytmów}
Zestawienie czasu i liczby porównań dla dwóch algorytmów zawiera
\cref{tab:porownanie-algorytmow}. Są to dane dydaktyczne, a nie wyniki pomiarów.

\begin{table}[htbp]
  \centering
  \caption{Porównanie przykładowych wyników dwóch algorytmów sortowania}
  \label{tab:porownanie-algorytmow}
  \begin{tabularx}{\textwidth}{@{}c *{4}{>{\centering\arraybackslash}X}@{}}
    \toprule
    Liczba elementów $n$
      & Bubble Sort czas [ms]
      & Insertion Sort czas [ms]
      & Bubble Sort porównania
      & Insertion Sort porównania \\
    \midrule
    100  & \num{0.4}   & \num{0.2}   & \num{4950}     & \num{2600}    \\
    500  & \num{8.9}   & \num{4.1}   & \num{124750}   & \num{65000}   \\
    1000 & \num{35.2}  & \num{16.3}  & \num{499500}   & \num{260000}  \\
    2000 & \num{141.5} & \num{65.1}  & \num{1999000}  & \num{1040000} \\
    5000 & \num{883.7} & \num{406.4} & \num{12497500} & \num{6500000} \\
    \bottomrule
  \end{tabularx}
\end{table}

Zależność czasu sortowania bąbelkowego od liczby elementów pokazuje
\cref{fig:wykres-sortowania}. Wykres dotyczy wyłącznie Bubble Sort.

\begin{figure}[htbp]
  \centering
  \includegraphics[width=0.8\textwidth]{%
    chapters/05-weryfikacja-i-wyniki/figures/wykres_sortowania.png}
  \caption{Przykładowa zależność czasu sortowania od liczby elementów}
  \label{fig:wykres-sortowania}
\end{figure}

\section{Ocena spełnienia wymagań}
Dla każdego wymagania podaj wynik testu oraz status: spełnione, częściowo spełnione albo niespełnione. Odnieś się do kryteriów z \cref{chap:cel-wymagania}.

\section{Dyskusja}
Zinterpretuj rezultaty, wyjaśnij odstępstwa i wskaż ograniczenia wiarygodności oceny.
